\input{uni}

\title{Отчёт по дисциплине «Нейронные сети и глубокое обучение»}
\author{Выполнил студент группы М-РИТ-25-1, Станиславчук Сергей}

\begin{document}

\setstretch{1.5}
\noindent
\begin{minipage}[t]{0.01\textwidth}
    \includegraphics{~/programmer/++Programmer/5_1/LSTU.jpg}
\end{minipage}%
\begin{minipage}[c]{1.1\linewidth}
    \centering
    {\small
    МИНИСТЕРСТВО НАУКИ И ВЫСШЕГО ОБРАЗОВАНИЯ РФ \\[1em]
    ФЕДЕРАЛЬНОЕ ГОСУДАРСТВЕННОЕ БЮДЖЕТНОЕ \\
    ОБРАЗОВАТЕЛЬНОЕ УЧРЕЖДЕНИЕ ВЫСШЕГО ОБРАЗОВАНИЯ \\
    «ЛИПЕЦКИЙ ГОСУДАРСТВЕННЫЙ ТЕХНИЧЕСКИЙ УНИВЕРСИТЕТ»
    }
\end{minipage}%

\vspace{3em}

\noindent
\hspace{1em}%
\begin{minipage}[t]{0.8\textwidth}
\textbf{Институт} \hspace{2em} компьютерных наук \\
\textbf{Кафедра} \hspace{2em} автоматизированных систем управления
\end{minipage}


\vspace{6em}

\begin{center}
\textbf{Курсовая работа} \\
по дисциплине \\ <<Нейронные сети и глубокое обучение>>
\end{center}

\vspace{6em}

\noindent
\begin{tabular*}{\textwidth}{@{\extracolsep{\fill}} l c r }
    Студент М-РИТ-25-1 & \makebox[4cm]{\hrulefill} & Станиславчук С.М. \\[2em]
Руководитель & \\
Профессор, д. тех. наук & \makebox[4cm]{\hrulefill} & Седых И.А. \\
\end{tabular*}

\vfill
\begin{center}
    Липецк, 2026
\end{center}

\thispagestyle{empty}

\newpage
\tableofcontents
\newpage

\section{Цель работы}

\begin{sloppypar}
    Исследование эффективности замены стандартных свёрток на depthwise separable convolutions в архитектуре ResNet-18. Сравнение двух вариантов модели — стандартного и с depthwise separable свёртками — по точности классификации, времени инференса на CPU и размеру модели. Обучение и оценка проводятся на датасете CIFAR-10.
\end{sloppypar}

\section{Задание}
\begin{sloppypar}
    Реализовать ResNet-18 для классификации изображений CIFAR-10 в двух вариантах: со~стандартными свёрточными слоями и с~depthwise separable свёртками (в~стиле MobileNet). Обучить обе модели. Измерить и сравнить следующие метрики: точность на тестовой выборке (Test Accuracy,~\%), латентность инференса на~CPU~(Latency,~мс), количество параметров~(Params) и~сериализованный размер модели~(Size,~МБ). Сделать вывод о~компромиссе между точностью и~эффективностью.
\end{sloppypar}

\newpage
\section{Исходный датасет}

Для обучения и тестирования используется датасет \textbf{CIFAR-10}~\cite{cifar10} — один из стандартных бенчмарков в задачах классификации изображений. Датасет содержит $60\,000$ цветных изображений размера $32 \times 32$ пикселя, равномерно распределённых по $10$ классам: самолёт, автомобиль, птица, кот, олень, собака, лягушка, лошадь, корабль, грузовик.

\begin{itemize}
    \item \textbf{Обучающая выборка:} $50\,000$ изображений (по $5\,000$ на класс).
    \item \textbf{Тестовая выборка:} $10\,000$ изображений (по $1\,000$ на класс).
    \item \textbf{Формат:} трёхканальные изображения RGB, значения пикселей в диапазоне $[0, 255]$.
\end{itemize}

\subsection{Предобработка данных}

\begin{enumerate}
    \item \textbf{Нормализация:} значения пикселей приводятся к диапазону $[0, 1]$ делением на $255$, затем нормализуются по каналам со средними $(0.4914, 0.4822, 0.4465)$ и стандартными отклонениями $(0.2470, 0.2435, 0.2616)$ — значения, вычисленные по обучающей выборке CIFAR-10.
    \item \textbf{Аугментация (только для обучения):}
    \begin{itemize}
        \item RandomCrop: случайная обрезка $32 \times 32$ с паддингом 4 пикселя.
        \item RandomHorizontalFlip: случайное отражение по горизонтали с вероятностью $0.5$.
    \end{itemize}
    \item \textbf{Для тестовой выборки} аугментация не применяется — только нормализация.
\end{enumerate}

Загрузка данных осуществляется через библиотеку \texttt{torchvision.datasets.CIFAR10} с автоматической загрузкой при первом обращении. Данные подаются в модель батчами размера $128$.

\section{Программа для обучения модели}

\subsection{Архитектура ResNet-18}

В качестве базовой архитектуры используется ResNet-18~\cite{resnet},
адаптированный для входного разрешения $32 \times 32$ (CIFAR-10). Стандартная
реализация ResNet-18 для ImageNet ($224 \times 224$) использует начальную
свёртку $7 \times 7$ со страйдом 2 и слой MaxPool, что для CIFAR-10 привело бы
к чрезмерно агрессивному уменьшению пространственного разрешения. Поэтому в
адаптированной версии:

\begin{itemize}
    \item Начальная свёртка заменена на $3 \times 3$ со страйдом 1 и паддингом 1.
    \item Слой MaxPool после начальной свёртки заменён на тождественное отображение (\texttt{nn.Identity}).
\end{itemize}

Архитектура состоит из следующих компонентов:

\begin{enumerate}
    \item \textbf{Начальный блок:} Conv2d($3 \to 64$, $3 \times 3$, stride~=~1, padding~=~1) $\to$ BatchNorm2d $\to$ ReLU.
    \item \textbf{Четыре Residual-блока} (\texttt{layer1}--\texttt{layer4}) с конфигурацией каналов $[64, 128, 256, 512]$ и числом BasicBlock-ов $[2, 2, 2, 2]$ соответственно. В~блоках \texttt{layer2}--\texttt{layer4} применяется страйд 2 для понижения пространственного разрешения.
    \item \textbf{BasicBlock} (рис.~\ref{fig:basicblock}): два свёрточных слоя $3\times3$ с BatchNorm и ReLU между ними, а также skip-connection (остаточная связь), прибавляющая вход блока к его выходу. Если размерности входа и выхода не совпадают (из-за изменения числа каналов или страйда), к skip-connection применяется понижающая свёртка $1 \times 1$ с BatchNorm.
    \item \textbf{Классификатор:} AdaptiveAvgPool2d($1 \times 1$) $\to$ Flatten $\to$ Linear($512 \to 10$).
\end{enumerate}

\begin{figure}[h]
    \centering
    \fbox{\parbox{0.85\textwidth}{
        \textbf{BasicBlock} \\
        Вход $x$ \\
        $\to$ Conv2d($C_{in} \to C_{out}$, $3\times3$) $\to$ BatchNorm $\to$ ReLU \\
        $\to$ Conv2d($C_{out} \to C_{out}$, $3\times3$) $\to$ BatchNorm \\
        $\to$ сложение с skip-connection($x$) $\to$ ReLU $\to$ выход
    }}
    \caption{Структура BasicBlock — базового остаточного блока ResNet-18.}
    \label{fig:basicblock}
\end{figure}

\subsection{Depthwise Separable Convolution}

В модифицированной версии модели все стандартные свёртки $3 \times 3$ и $1 \times 1$ заменены на depthwise separable convolution~\cite{mobilenet} — двухэтапную операцию:

\begin{enumerate}
    \item \textbf{Depthwise convolution:} свёртка $k \times k$, применяемая независимо к каждому входному каналу (\texttt{groups}~=~\texttt{in\_channels}). Количество выходных каналов равно количеству входных. Не смешивает информацию между каналами.
    \item \textbf{Pointwise convolution:} свёртка $1 \times 1$, смешивающая каналы и изменяющая их количество до требуемого выходного.
\end{enumerate}

Такая замена радикально сокращает число параметров и вычислительную сложность свёрточных слоёв. Для свёртки с $C_{in}$ входными каналами, $C_{out}$ выходными и ядром $k \times k$:

\begin{itemize}
    \item \textbf{Стандартная:} $k^2 \cdot C_{in} \cdot C_{out}$ параметров.
    \item \textbf{Depthwise separable:} $k^2 \cdot C_{in} + C_{in} \cdot C_{out}$ параметров.
\end{itemize}

При $k = 3$, $C_{in} = C_{out} = 64$ сокращение составляет примерно $\frac{9 \cdot 64 \cdot 64}{9 \cdot 64 + 64 \cdot 64} \approx 7.9$ раз для данного слоя.

\subsection{Процесс обучения}

Обучение обеих моделей выполняется в идентичных условиях. Гиперпараметры:

\begin{itemize}
    \item \textbf{Оптимизатор:} SGD с моментом $0.9$ и weight decay $5 \cdot 10^{-4}$.
    \item \textbf{Начальная скорость обучения (lr):} $0.1$.
    \item \textbf{Планировщик lr:} MultiStepLR с уменьшением в $10$ раз на эпохах, соответствующих $50\%$ и $75\%$ общего числа эпох (для $20$ эпох — milestones $[10, 15]$, $\gamma = 0.1$).
    \item \textbf{Функция потерь:} CrossEntropyLoss.
    \item \textbf{Число эпох:} $20$.
    \item \textbf{Размер батча:} $128$.
    \item \textbf{Инициализация весов:} Kaiming Normal (He-инициализация) для свёрток, константная ($1.0$ для BatchNorm-весов, $0.0$ для смещений).
\end{itemize}

На каждой эпохе модель обучается на полной тренировочной выборке, после чего вычисляется точность на тестовой выборке. Сохраняется чекпоинт с наилучшей точностью. Для обеих моделей используется одно и то же начальное значение random seed (42).

\subsection{Методика бенчмарка}

После обучения для каждой модели выполняются замеры на CPU:

\begin{enumerate}
    \item \textbf{Количество параметров:} подсчитывается сумма \texttt{numel()} всех тензоров с \texttt{requires\_grad=True}.
    \item \textbf{Размер модели:} сериализация \texttt{state\_dict()} в буфер через \texttt{torch.save}, размер буфера в мегабайтах.
    \item \textbf{Латентность инференса:} измеряется время одного прямого прохода (batch size = 1, вход $3\times32\times32$) на CPU. Выполняется $50$ прогревочных прогонов, затем $300$ замеров. Результат — среднее арифметическое $\pm$ стандартное отклонение в миллисекундах.
\end{enumerate}

\subsection{Результаты экспериментов}

Обучение проводилось на GPU (NVIDIA), бенчмарк — на CPU с фиксированным числом потоков (\texttt{OMP\_NUM\_THREADS=1}, \texttt{torch.set\_num\_threads(1)}). Результаты приведены в таблице~\ref{tab:results}.

\begin{table}[h]
    \centering
    \caption{Сравнение стандартного и depthwise separable ResNet-18.}
    \label{tab:results}
    \begin{tabular}{|l|c|c|}
        \hline
        \textbf{Метрика} & \textbf{Standard} & \textbf{Separable} \\
        \hline
        Параметры & $11\,173\,962$ & $1\,438\,565$ \\
        Размер (МБ) & $42.70$ & $5.57$ \\
        Latency CPU (мс) & $27.53 \pm 2.33$ & $6.93 \pm 0.16$ \\
        Test Accuracy (\%) & $91.09$ & $89.81$ \\
        \hline
    \end{tabular}
\end{table}

\begin{table}[h]
    \centering
    \caption{Относительные показатели separable-модели.}
    \label{tab:relative}
    \begin{tabular}{|l|c|}
        \hline
        \textbf{Показатель} & \textbf{Значение} \\
        \hline
        Ускорение инференса & $\times 3.97$ \\
        Сокращение размера & $\times 7.67$ \\
        Падение точности & $+1.28$ п.п. \\
        \hline
    \end{tabular}
\end{table}

\section{Вывод}

Проведённое исследование демонстрирует, что замена стандартных свёрток на depthwise separable в архитектуре ResNet-18 позволяет:

\begin{enumerate}
    \item \textbf{Сократить число параметров} в $7.67$ раз (с $11.17$ млн до $1.44$ млн) и размер сериализованной модели с $42.70$ МБ до $5.57$ МБ.
    \item \textbf{Ускорить инференс на CPU} в $3.97$ раза (с $27.53$ мс до $6.93$ мс на одно изображение).
    \item \textbf{Сохранить точность} на уровне $89.81\%$ против $91.09\%$ у стандартной модели — потеря составляет всего $1.28$ процентных пункта.
\end{enumerate}

Полученные результаты согласуются с выводами оригинальной работы по MobileNet~\cite{mobilenet}: depthwise separable convolutions обеспечивают значительный выигрыш в эффективности при незначительном снижении точности. Это делает их привлекательным выбором для развёртывания CNN на устройствах с ограниченными вычислительными ресурсами (CPU, мобильные и встраиваемые платформы).

Наблюдаемый скачок точности на 10-й эпохе (с $79.5\%$ до $88.9\%$ для стандартной модели) объясняется срабатыванием планировщика MultiStepLR — уменьшение learning rate в $10$ раз позволяет оптимизатору совершить более точный спуск и выйти из локальных минимумов функции потерь.

\newpage
\section{Список источников}

\begin{enumerate}
    \begin{sloppypar}
        \item He~K., Zhang~X., Ren~S., Sun~J. Deep Residual Learning for Image Recognition //
            Proceedings of the IEEE Conference on Computer Vision and Pattern Recognition
            (CVPR). — 2016. — P.~770--778.
            \label{resnet}

        \item Howard~A.G. et~al. MobileNets: Efficient Convolutional Neural Networks for
            Mobile Vision Applications // arXiv preprint arXiv:1704.04861. — 2017.
            \label{mobilenet}

        \item Krizhevsky~A. Learning Multiple Layers of Features from Tiny Images //
            Technical Report, University of Toronto. — 2009.
            \label{cifar10}

        \item CIFAR-10 — Canadian Institute for Advanced Research [Электронный ресурс] /
            Режим доступа к данным: https://www.cs.toronto.edu/~kriz/cifar.html,
            свободный.

        \item PyTorch Documentation: torchvision.datasets.CIFAR10 [Электронный ресурс] /
            Режим доступа к данным: https://pytorch.org/vision/stable/generated/
            torchvision.datasets.CIFAR10.html, свободный.

        \item Chollet~F. Xception: Deep Learning with Depthwise Separable Convolutions //
            Proceedings of the IEEE Conference on Computer Vision and Pattern Recognition
            (CVPR). — 2017. — P.~1251--1258.
    \end{sloppypar}
\end{enumerate}

\newpage
\section*{Приложение 1. Код программы}

\footnotesize
\setlength{\baselineskip}{0.25\baselineskip}

\subsection*{models.py — Архитектура ResNet-18 и DepthwiseSeparableConv2d}

\begin{verbatim}
"""ResNet-18 for CIFAR-10: standard and depthwise-separable convolutions."""

from __future__ import annotations

import io
from typing import Type

import torch
import torch.nn as nn


class DepthwiseSeparableConv2d(nn.Module):
    """MobileNet-style block: depthwise kxk + pointwise 1x1."""

    def __init__(
        self,
        in_channels: int,
        out_channels: int,
        kernel_size: int,
        stride: int = 1,
        padding: int = 0,
        bias: bool = False,
    ) -> None:
        super().__init__()
        self.depthwise = nn.Conv2d(
            in_channels,
            in_channels,
            kernel_size,
            stride=stride,
            padding=padding,
            groups=in_channels,
            bias=bias,
        )
        self.pointwise = nn.Conv2d(in_channels, out_channels, 1, bias=bias)

    def forward(self, x: torch.Tensor) -> torch.Tensor:
        return self.pointwise(self.depthwise(x))


def conv3x3(
    in_channels: int,
    out_channels: int,
    stride: int = 1,
    *,
    separable: bool = False,
) -> nn.Module:
    if separable:
        return DepthwiseSeparableConv2d(
            in_channels, out_channels, 3, stride=stride, padding=1, bias=False
        )
    return nn.Conv2d(
        in_channels, out_channels, 3, stride=stride, padding=1, bias=False
    )


class BasicBlock(nn.Module):
    expansion = 1

    def __init__(
        self,
        in_channels: int,
        out_channels: int,
        stride: int = 1,
        downsample: nn.Module | None = None,
        *,
        separable: bool = False,
    ) -> None:
        super().__init__()
        self.conv1 = conv3x3(in_channels, out_channels, stride, separable=separable)
        self.bn1 = nn.BatchNorm2d(out_channels)
        self.conv2 = conv3x3(out_channels, out_channels, separable=separable)
        self.bn2 = nn.BatchNorm2d(out_channels)
        self.relu = nn.ReLU(inplace=True)
        self.downsample = downsample

    def forward(self, x: torch.Tensor) -> torch.Tensor:
        identity = x
        out = self.relu(self.bn1(self.conv1(x)))
        out = self.bn2(self.conv2(out))
        if self.downsample is not None:
            identity = self.downsample(x)
        return self.relu(out + identity)


class ResNet(nn.Module):
    def __init__(
        self,
        block: Type[BasicBlock],
        layers: list[int],
        num_classes: int = 10,
        *,
        separable: bool = False,
    ) -> None:
        super().__init__()
        self.in_channels = 64
        self.separable = separable

        self.conv1 = conv3x3(3, 64, separable=separable)
        self.bn1 = nn.BatchNorm2d(64)
        self.relu = nn.ReLU(inplace=True)
        self.maxpool = nn.Identity()

        self.layer1 = self._make_layer(block, 64, layers[0])
        self.layer2 = self._make_layer(block, 128, layers[1], stride=2)
        self.layer3 = self._make_layer(block, 256, layers[2], stride=2)
        self.layer4 = self._make_layer(block, 512, layers[3], stride=2)

        self.avgpool = nn.AdaptiveAvgPool2d((1, 1))
        self.fc = nn.Linear(512 * block.expansion, num_classes)
        self._init_weights()

    def _make_layer(
        self,
        block: Type[BasicBlock],
        out_channels: int,
        blocks: int,
        stride: int = 1,
    ) -> nn.Sequential:
        downsample = None
        if stride != 1 or self.in_channels != out_channels:
            downsample = nn.Sequential(
                nn.Conv2d(self.in_channels, out_channels, 1, stride=stride, bias=False),
                nn.BatchNorm2d(out_channels),
            )
        layers = [
            block(self.in_channels, out_channels, stride, downsample,
                  separable=self.separable)
        ]
        self.in_channels = out_channels
        for _ in range(1, blocks):
            layers.append(
                block(self.in_channels, out_channels, separable=self.separable)
            )
        return nn.Sequential(*layers)

    def _init_weights(self) -> None:
        for m in self.modules():
            if isinstance(m, nn.Conv2d):
                nn.init.kaiming_normal_(m.weight, mode="fan_out",
                                        nonlinearity="relu")
            elif isinstance(m, nn.BatchNorm2d):
                nn.init.constant_(m.weight, 1)
                nn.init.constant_(m.bias, 0)

    def forward(self, x: torch.Tensor) -> torch.Tensor:
        x = self.relu(self.bn1(self.conv1(x)))
        x = self.maxpool(x)
        x = self.layer4(self.layer3(self.layer2(self.layer1(x))))
        x = self.avgpool(x)
        x = torch.flatten(x, 1)
        return self.fc(x)


def resnet18(num_classes: int = 10, *, separable: bool = False) -> ResNet:
    return ResNet(BasicBlock, [2, 2, 2, 2], num_classes=num_classes,
                  separable=separable)


def count_parameters(model: nn.Module) -> int:
    return sum(p.numel() for p in model.parameters() if p.requires_grad)


def model_size_bytes(model: nn.Module) -> int:
    buffer = io.BytesIO()
    torch.save(model.state_dict(), buffer)
    return buffer.tell()
\end{verbatim}

\subsection*{train.py — Обучение и оценка}

\begin{verbatim}
"""Training ResNet-18 on CIFAR-10."""

from __future__ import annotations

import time
from pathlib import Path

import torch
import torch.nn as nn
from torch.utils.data import DataLoader
from torchvision import datasets, transforms

from models import ResNet, resnet18


def get_dataloaders(
    data_dir: str | Path,
    batch_size: int = 128,
    num_workers: int = 2,
) -> tuple[DataLoader, DataLoader]:
    data_dir = Path(data_dir)
    normalize = transforms.Normalize(
        (0.4914, 0.4822, 0.4465),
        (0.2470, 0.2435, 0.2616),
    )
    train_tf = transforms.Compose([
        transforms.RandomCrop(32, padding=4),
        transforms.RandomHorizontalFlip(),
        transforms.ToTensor(),
        normalize,
    ])
    test_tf = transforms.Compose([transforms.ToTensor(), normalize])

    train_set = datasets.CIFAR10(
        root=data_dir, train=True, download=True, transform=train_tf
    )
    test_set = datasets.CIFAR10(
        root=data_dir, train=False, download=True, transform=test_tf
    )

    train_loader = DataLoader(
        train_set, batch_size=batch_size, shuffle=True,
        num_workers=num_workers, pin_memory=torch.cuda.is_available(),
    )
    test_loader = DataLoader(
        test_set, batch_size=batch_size, shuffle=False,
        num_workers=num_workers, pin_memory=torch.cuda.is_available(),
    )
    return train_loader, test_loader


@torch.no_grad()
def evaluate(model: nn.Module, loader: DataLoader,
             device: torch.device) -> float:
    model.eval()
    correct = 0
    total = 0
    for images, labels in loader:
        images, labels = images.to(device), labels.to(device)
        logits = model(images)
        preds = logits.argmax(dim=1)
        correct += (preds == labels).sum().item()
        total += labels.size(0)
    return 100.0 * correct / total


def train_model(
    model: ResNet,
    train_loader: DataLoader,
    test_loader: DataLoader,
    *,
    epochs: int = 50,
    lr: float = 0.1,
    device: torch.device | None = None,
    save_path: Path | None = None,
) -> dict[str, float]:
    device = device or torch.device(
        "cuda" if torch.cuda.is_available() else "cpu"
    )
    model = model.to(device)

    criterion = nn.CrossEntropyLoss()
    optimizer = torch.optim.SGD(
        model.parameters(), lr=lr, momentum=0.9, weight_decay=5e-4
    )
    scheduler = torch.optim.lr_scheduler.MultiStepLR(
        optimizer, milestones=[int(epochs * 0.5), int(epochs * 0.75)],
        gamma=0.1
    )

    best_acc = 0.0
    t0 = time.perf_counter()

    for epoch in range(1, epochs + 1):
        model.train()
        running_loss = 0.0
        for images, labels in train_loader:
            images, labels = images.to(device), labels.to(device)
            optimizer.zero_grad()
            loss = criterion(model(images), labels)
            loss.backward()
            optimizer.step()
            running_loss += loss.item()

        scheduler.step()
        train_loss = running_loss / len(train_loader)
        test_acc = evaluate(model, test_loader, device)
        best_acc = max(best_acc, test_acc)
        print(f"Epoch {epoch:3d}/{epochs}  loss={train_loss:.4f}  "
              f"test_acc={test_acc:.2f}%  best={best_acc:.2f}%")

    elapsed = time.perf_counter() - t0
    final_acc = evaluate(model, test_loader, device)

    if save_path is not None:
        save_path.parent.mkdir(parents=True, exist_ok=True)
        torch.save({
            "state_dict": model.state_dict(),
            "separable": model.separable,
            "test_acc": final_acc,
            "best_acc": best_acc,
        }, save_path)

    return {"final_acc": final_acc, "best_acc": best_acc,
            "train_time_s": elapsed}
\end{verbatim}

\subsection*{benchmark.py — Замер латентности и размера модели}

\begin{verbatim}
"""CPU latency and model size comparison."""

from __future__ import annotations

import time
from dataclasses import dataclass

import torch
import torch.nn as nn

from models import count_parameters, model_size_bytes


@dataclass
class BenchmarkResult:
    name: str
    params: int
    size_mb: float
    latency_ms: float
    latency_std_ms: float
    test_acc: float | None = None


@torch.no_grad()
def measure_latency(
    model: nn.Module,
    *,
    input_shape: tuple[int, ...] = (1, 3, 32, 32),
    warmup: int = 20,
    repeats: int = 100,
    device: torch.device | None = None,
) -> tuple[float, float]:
    device = device or torch.device("cpu")
    model = model.to(device)
    model.eval()

    x = torch.randn(*input_shape, device=device)
    for _ in range(warmup):
        model(x)

    times: list[float] = []
    for _ in range(repeats):
        t0 = time.perf_counter()
        model(x)
        times.append((time.perf_counter() - t0) * 1000.0)

    mean = sum(times) / len(times)
    variance = sum((t - mean) ** 2 for t in times) / len(times)
    return mean, variance ** 0.5


def benchmark_model(
    model: nn.Module,
    name: str,
    *,
    test_acc: float | None = None,
    device: torch.device | None = None,
    warmup: int = 20,
    repeats: int = 100,
) -> BenchmarkResult:
    device = device or torch.device("cpu")
    model_cpu = model.to(device)
    model_cpu.eval()

    lat_mean, lat_std = measure_latency(
        model_cpu, warmup=warmup, repeats=repeats, device=device,
    )
    return BenchmarkResult(
        name=name,
        params=count_parameters(model_cpu),
        size_mb=model_size_bytes(model_cpu) / (1024 * 1024),
        latency_ms=lat_mean,
        latency_std_ms=lat_std,
        test_acc=test_acc,
    )
\end{verbatim}

\end{document}
